\documentclass{slemanslides}

\title{GPIO in Practice}
\course{Learning Embedded Systems}
\lesson{Lesson 03}

\begin{document}

\cover

\divider{What is a GPIO}

\begin{frame}{Concepts}
  \concept{GPIO}{General Purpose Input/Output: a microcontroller pin
  that firmware configures as a digital input or output.}

  \concept{Register}{A memory-mapped address. Writing to it changes the
  physical state of a peripheral.}
\end{frame}



\begin{frame}{Pin modes}
  \begin{itemize}
    \item \term{Input:} reads a logic level (button, sensor)
      \begin{itemize}
        \item pull-up, pull-down or floating
      \end{itemize}
    \item \term{Output:} drives a load (LED, relay, transistor)
      \begin{itemize}
        \item push-pull or open-drain
      \end{itemize}
    \item \term{Alternate function:} UART, SPI, I2C, PWM
    \item \term{Analog:} ADC input
  \end{itemize}
\end{frame}




\begin{frame}{Turning on an LED}
  \begin{enumerate}
    \item Enable the port clock (\cc{RCC->AHB1ENR})
    \item Set the pin as output (\cc{GPIOx->MODER})
    \item Write to the output register (\cc{GPIOx->ODR})
  \end{enumerate}

  \warn{Without the clock enabled, writes to the port registers are
  silently ignored.}
\end{frame}



\begin{frame}{Image with text}
  \begin{columns}[c]
    \begin{column}{0.55\textwidth}
      \includegraphics[width=\linewidth]{sample}
    \end{column}
    \begin{column}{0.4\textwidth}
      \begin{itemize}
        \item Green LED on \cc{PA5}
        \item Blue button on \cc{PC13}
        \item On-board ST-LINK
      \end{itemize}
    \end{column}
  \end{columns}
  \source{NUCLEO-F401RE user manual, ST.}
\end{frame}



\begin{frame}{Centered figure}
  \fig{sample}{GPIO block diagram.}
\end{frame}



%% \fullscreen{file} or \fullscreen[caption]{file}
\fullscreen[NUCLEO-F401RE board]{sample}



\begin{frame}{Three images}
  \begin{columns}[t]
    \begin{column}{0.31\textwidth}\centering
      \includegraphics[width=\linewidth]{sample}\\[3pt]
      {\footnotesize\color{muted}Push-pull}
    \end{column}
    \begin{column}{0.31\textwidth}\centering
      \includegraphics[width=\linewidth]{sample}\\[3pt]
      {\footnotesize\color{muted}Open-drain}
    \end{column}
    \begin{column}{0.31\textwidth}\centering
      \includegraphics[width=\linewidth]{sample}\\[3pt]
      {\footnotesize\color{muted}Input with pull-up}
    \end{column}
  \end{columns}
\end{frame}



\divider{Hands on code}



%% Code (frame needs [fragile])
%% \begin{code}[listings options]{file name}   (empty {} = no bar)
%% Fits ~11 lines at the default size; for longer snippets pass
%% [basicstyle=\ttfamily\footnotesize\color{codefg}].
\begin{frame}[fragile]{Bare-metal blink}
\begin{code}{main.c}
#include "stm32f4xx.h"

int main(void)
{
    RCC->AHB1ENR |= RCC_AHB1ENR_GPIOAEN;    // enable port A clock
    GPIOA->MODER |= (1U << (5 * 2));        // PA5 as output
    while (1) {
        GPIOA->ODR ^= (1U << 5);            // toggle the LED
        for (volatile uint32_t i = 0; i < 500000; i++);
    }
}
\end{code}
\end{frame}



\begin{frame}[fragile]{Code and explanation}
  \begin{columns}[t]
    \begin{column}{0.56\textwidth}
\begin{code}{gpio.h}
typedef struct {
    volatile uint32_t MODER;
    volatile uint32_t OTYPER;
    volatile uint32_t OSPEEDR;
    volatile uint32_t PUPDR;
    volatile uint32_t IDR;
    volatile uint32_t ODR;
} GPIO_TypeDef;
\end{code}
    \end{column}
    \begin{column}{0.4\textwidth}
      \begin{itemize}
        \item Each field is a 32-bit register
        \item Field order follows the \term{offsets} in the reference
          manual
        \item \cc{volatile}: the compiler cannot optimize the access away
      \end{itemize}
    \end{column}
  \end{columns}
\end{frame}




\begin{frame}[fragile]{Build and flash}
\begin{code}[language=bash,numbers=none]{terminal}
$ make
$ st-flash write build/blink.bin 0x08000000
\end{code}

  \tip{Use \cc{numbers=none} for terminal commands.}
\end{frame}




\begin{frame}{Callouts}
  \info{Neutral note --- \texttt{\textbackslash info\{...\}}}
  \tip{Shortcut or good practice --- \texttt{\textbackslash tip\{...\}}}
  \warn{Something to watch out for --- \texttt{\textbackslash warn\{...\}}}
  \danger{Something that fries the board --- \texttt{\textbackslash danger\{...\}}}
\end{frame}



\begin{frame}{Text highlight}
  \blue{blue}, \orange{orange}, \red{red}, \green{green},
  \term{term} and \cc{inline_code()}.

  \medskip
  Works with \textbf{bold} and \uline{underline}.
\end{frame}



\closing

\end{document}